\title{The Era of 1-bit LLMs: All Large Language Models are in 1.58 Bits}

\begin{abstract}
Recent advancements, exemplified by BitNet~\cite{bitnet}, are ushering in a new generation of 1-bit Large Language Models (LLMs). Here, we present a new 1-bit LLM model, designated as \textbf{\text{BitNet b1.58}}, wherein each parameter of the model adopts a ternary value from the set \{-1, 0, 1\}. This approach achieves performance parity with full-precision Transformer LLMs—whether FP16 or BF16—in terms of perplexity and end-task results, provided equivalent model scale and training data. Importantly, this design offers substantial improvements in latency, memory usage, throughput, and energy efficiency. Notably, the 1.58-bit LLM formulation establishes a new paradigm for scaling laws and methodologies in LLM training, balancing both efficiency and performance. Additionally, it points to innovative avenues for computation, fostering the development of hardware specialized for 1-bit LLM operations.
\end{abstract}

\section{The Era of 1-bit LLMs}

The expansion of Large Language Models (LLMs) within artificial intelligence has been dramatic in both scale and performance, achieving state-of-the-art results in diverse natural language processing applications. Nonetheless, this upscaling brings practical deployment obstacles and accentuates environmental and financial concerns due to heightened energy requirements.

To mitigate these issues, researchers have explored post-training quantization as a strategy for developing inference-efficient, low-bit models~\cite{smoothquant, gptq, quip, quip_sharp}. This method compresses both the weights and activations by lowering their numerical precision, thereby drastically decreasing the resource demands of LLMs. While the move from 16-bit representations down to 4-bit models has garnered momentum~\cite{gptq, awq}, such post-training quantization, though prevalent in industry-standard LLMs, does not yield optimal outcomes.

Recent advances in 1-bit model architectures, exemplified by BitNet~\cite{bitnet}, have shown potential for substantially lowering the operational costs of LLMs while preserving model accuracy. Traditional LLMs utilize 16-bit floating point representations (either FP16 or BF16), and their computational core largely consists of matrix multiplications. This computation is dominated by the expense of floating-point addition and multiplication. Conversely, BitNet architectures perform matrix multiplications using only integer additions, which results in orders-of-magnitude reductions in energy consumption. As computational performance on many hardware accelerators is often constrained by power limits, such reductions in energy use also enable faster model execution.

Beyond raw computation, moving model weights from off-chip DRAM to on-chip accelerator memory (such as SRAM) imposes significant overhead during inference. While increasing on-chip SRAM capacity has been explored as a strategy to boost throughput, it comes with much higher expense relative to DRAM. Owing to their compressed parameter representations, 1-bit LLMs offer greatly diminished demands on both memory capacity and bandwidth when compared to full-precision counterparts. This translates into more efficient and rapid loading of weights from DRAM, yielding improvements in inference speed and cost-effectiveness.

Here, we present a notable extension to the 1-bit LLM design, introducing \textbf{\text{BitNet b1.58}}, in which every parameter is ternary, assuming values from the set \{-1, 0, 1\}. By augmenting the original 1-bit BitNet with an additional zero value, the representation expands to 1.58 bits per parameter. \text{BitNet b1.58} maintains all the core advantages of its 1-bit predecessor, notably a computation regime that eschews most multiplication operations in matrix multiplications and is thus highly amenable to efficient optimization. Its memory usage, throughput, and inference latency are all superior relative to FP16 LLMs, while energy demands remain comparable to the original 1-bit BitNet. Moreover, \text{BitNet b1.58} offers two key improvements: (1) It enhances modeling power through explicit feature filtering, enabled by allowing model weights to be zero, which substantially bolsters the efficacy of 1-bit LLMs; (2) Our empirical studies indicate that, beginning at model sizes of 3B parameters, \text{BitNet b1.58} can achieve parity with full-precision (FP16) baselines regarding perplexity and downstream task performance, under identical settings (such as model scale and training data volume).

\section{BitNet b1.58}

\text{BitNet b1.58} builds upon the BitNet framework, a Transformer variant that substitutes the standard \emph{nn.Linear} layers with \emph{BitLinear} modules. Training is conducted from the ground up, utilizing weights quantized to 1.58 bits and activations quantized to 8 bits. This version incorporates a number of updates beyond the initial BitNet design, which are outlined below.

\paragraph{Quantization Function.} To ensure weight values are limited to -1, 0, or +1, we employ the \emph{absmean} quantization approach. This method normalizes the weight matrix using its mean absolute value, after which each element is rounded to the closest of the set \{-1, 0, +1\}:

\begin{equation}
    \widetilde{W} = \mathrm{RoundClip}(\frac{W}{\gamma + \epsilon}, -1, 1),
\end{equation}

\begin{equation}
    \mathrm{RoundClip}(x, a, b) = \max(a, \min(b, \mathrm{round}(x))),
\end{equation}

\begin{equation}
    \gamma = \frac{1}{nm}\sum_{ij} |W_{ij}|.
\end{equation}

For the activation quantization, we largely retain the procedure from the original BitNet. However, rather than normalizing pre-activation values to the interval $[0, Q_b]$, activations are scaled within $[-Q_b, Q_b]$ for each token, thereby eliminating zero-point quantization. This modification simplifies both coding and system-level integration and, based on our empirical results, does not significantly affect model accuracy.

\paragraph{LLaMA-alike Components.}

With LLaMA~\cite{llama, llama2} having emerged as the foundational architecture for many open-source LLMs, our implementation of \text{BitNet b1.58} incorporates LLaMA-style components to align with open-source conventions. Namely, it employs RMSNorm~\cite{rmsnorm}, SwiGLU~\cite{swiglu}, rotary embedding~\cite{rope}, and omits bias terms. As a result, \text{BitNet b1.58} can be seamlessly adopted into mainstream open-source libraries (such as Huggingface, vLLM~\cite{vllm}, and llama.cpp\footnote{\url{https://github.com/ggerganov/llama.cpp}}), requiring only minimal adaptation.

\section{Results}

We conducted a comparative analysis between \text{BitNet b1.58} and our reproduced FP16 \text{LLaMA LLM} across multiple model scales. Both models underwent pre-training on the RedPajama dataset~\cite{redpajama}, each exposed to 100 billion tokens, to maintain an equivalent experimental setting. For evaluation, we assessed their zero-shot capabilities on several language benchmarks, including ARC-Easy~\cite{arc}, ARC-Challenge~\cite{arc}, Hellaswag~\cite{hellaswag}, Winogrande~\cite{winoGrande}, PIQA~\cite{piqa}, OpenbookQA~\cite{openbookqa}, and BoolQ~\cite{boolq}. Additionally, validation perplexity metrics were collected on WikiText2~\cite{wikitext2} and C4~\cite{c4} datasets.

We also investigated both the GPU memory footprint and inference latency during runtime for \text{LLaMA LLM} and \text{BitNet b1.58}. Measurements were taken using the FasterTransformer\footnote{\url{https://github.com/NVIDIA/FasterTransformer}} framework, which delivers highly efficient large language model inference on GPUs. The inference process for \text{BitNet b1.58} additionally leveraged the 2-bit kernel from Ladder~\cite{ladder}. Time per generated token was recorded, as this metric predominantly determines overall inference cost.

A summary of the models’ perplexity and resource usage is presented in Table~\ref{tab:ppl}. The findings indicate that \text{BitNet b1.58} achieves parity in perplexity with the FP16 \text{LLaMA LLM} beginning at the 3B parameter level, while offering a 2.71-fold speedup and utilizing only 28\% (3.55x reduction) of the GPU memory. Notably, the 3.9B variant of \text{BitNet b1.58} exhibits a 2.4x reduction in latency, requires 3.32 times less memory, and surpasses the performance of the 3B \text{LLaMA LLM}.

A detailed breakdown of zero-shot task accuracy is shown in Table~\ref{tab:zero-shot}, using the standard \emph{lm-evaluation-harness}\footnote{\url{https://github.com/EleutherAI/lm-evaluation-harness}} pipeline for assessment. The results demonstrate that the performance disparity between \text{BitNet b1.58} and \text{LLaMA LLM} diminishes with growing model scale. Crucially, \text{BitNet b1.58} attains parity with the full-precision baseline at the 3B parameter size. In line with the perplexity analysis, results on end tasks show that \text{BitNet b1.58} 3.9B exceeds the accuracy of \text{LLaMA LLM} 3B, with further reductions in both memory and latency demands. Collectively, these results indicate that \text{BitNet b1.58} offers a Pareto-optimal enhancement over leading LLM baselines.

\paragraph{Memory and Latency}

We expanded our evaluation to include models with sizes of 7B, 13B, and 70B parameters, systematically analyzing the associated computational costs. As depicted in Figure~\ref{fig:latency-memory-energy}, both latency and memory usage exhibit characteristic patterns as model size increases. Notably, speed-up is amplified for larger models: \text{BitNet b1.58} at 70B achieves a 4.1$\times$ acceleration over the \text{LLaMA LLM} reference. This improvement stems from the proportional increase in \emph{nn.Linear} computational demands as models grow. Similarly, memory usage displays parallel scaling behavior; because the embedding layers retain full-precision representation, their share of overall memory decreases in larger architectures. The reported latency and memory statistics rely on a 2-bit kernel implementation, indicating potential for further gains through optimization.

\paragraph{Energy}

In addition, we estimated the energy expenditure for arithmetic operations in both \text{BitNet b1.58} and \text{LLaMA LLM}. The primary contributor to LLM computational expense is matrix multiplication, and thus we center our analysis on this operation. As presented in Figure~\ref{fig:energy}, \text{BitNet b1.58} predominantly leverages INT8 addition operations, while \text{LLaMA LLM} employs a mix of FP16 addition and multiplication. Utilizing the energy consumption framework described in~\cite{energycost, pokebnn}, we find that \text{BitNet b1.58} reduces the arithmetic energy cost for matrix multiplication by a factor of 71.4$\times$ on 7nm process nodes. Additionally, our comprehensive end-to-end measurements—performed on models processing 512 tokens—reveal that as model capacity increases, \text{BitNet b1.58} demonstrates progressively superior energy efficiency compared to the FP16-based \text{LLaMA LLM}. This is primarily because the proportion of computation attributed to \emph{nn.Linear} expands with model size, whereas the contribution of other architectural components becomes comparatively marginal for larger models.

\paragraph{Throughput}

To evaluate throughput, we conducted a comparison between \text{BitNet b1.58} and \text{LLaMA LLM} (each containing 70B parameters) across two 80GB A100 GPUs, utilizing pipeline parallelism~\cite{gpipe} to enable \text{LLaMA LLM} 70B to operate within device memory constraints. We progressively increased the batch size up to the point of exhausting GPU memory, using a fixed sequence length of 512 tokens. The data summarized in Table~\ref{tab:throughput} demonstrates that \text{BitNet b1.58} 70B accommodates up to 11 times larger batch sizes than \text{LLaMA LLM}, yielding an 8.9-fold increase in throughput.

\textbf{\text{BitNet b1.58} establishes a novel scaling relationship linking model quality and inference efficiency.} Specifically, the following comparisons—grounded in the evidence shown in Figure~\ref{fig:latency-memory-energy} and \ref{fig:energy}—offer equivalences across model sizes for 1.58-bit and 16-bit variants:

\begin{itemize}
    \item In terms of latency, memory footprint, and power usage, 13B BitNet b1.58 surpasses the efficiency of a 3B FP16 LLM.
    \item Similarly, the 30B BitNet b1.58 exhibits greater efficiency by these same measures compared to a 7B FP16 LLM.
    \item The 70B BitNet b1.58 outperforms a 13B FP16 LLM in latency, memory consumption, and energy use.
\end{itemize}

\paragraph{Training with 2T Tokens}

Token count during training plays a pivotal role for LLMs. To assess how \text{BitNet b1.58} scales with increasing tokens, we trained it using 2T tokens in alignment with the StableLM-3B~\cite{StableLM-3B-4E1T} data procedure, which is employed by the leading open-source 3B model. Evaluations of both models were carried out on a composite benchmark, which includes Winogrande~\cite{winoGrande}, PIQA~\cite{piqa}, SciQ~\cite{sciq}, LAMBADA~\cite{lambada}, and ARC-easy~\cite{arc}. Zero-shot performance metrics are provided in Table~\ref{tab:2t}. For metrics involving both accuracy and normalized accuracy, their mean was used. Results for StableLM 3B at the 2T token scale are sourced from its published technical documentation. Our results reveal that \text{BitNet b1.58} consistently achieves higher performance across every evaluated task, supporting the claim that 1.58-bit LLMs are capable of strong generalization.

\section{Discussion and Future Work}

\textbf{1-bit Mixture-of-Experts (MoE) LLMs}

Mixture-of-Experts (MoE) architectures offer a resource-efficient methodology for large language models (LLMs) by considerably reducing computational FLOPs. Nevertheless, substantial memory usage and significant inter-chip communication demands have been barriers to their widespread adoption. The advent of 1.58-bit LLMs provides viable solutions to these issues. By decreasing the overall memory requirements, fewer devices are necessary for MoE deployment. This, in turn, sharply lowers the latency and cost associated with network-based activation transfer. Ideally, these reductions could ultimately eliminate such overhead altogether if the complete model is able to reside on a single chip.

\textbf{Native Support of Long Sequence in LLMs}

Managing long input sequences has become increasingly pivotal for LLM applications. A persistent difficulty in long-sequence inference is the extensive memory required to store key-value (KV) caches. \text{BitNet b1.58} moves the field forward by cutting activation precision from 16 bits to 8 bits, effectively allowing the same hardware to handle twice the sequence length. Even further, for 1.58-bit LLMs, activations may be compressed losslessly to 4 bits or less, a promising direction for future exploration.

\textbf{LLMs on Edge and Mobile}

Adopting 1.58-bit LLMs has the potential to transform language model deployment on edge and mobile hardware, which are commonly constrained by both memory capacity and processing power. The lower resource requirements of these quantized models make it feasible to bring sophisticated LLM capabilities to devices previously considered unsuitable. This expansion supports broader application development and enables mobile and edge devices to take advantage of LLM technologies in ways not previously feasible. In addition, 1.58-bit LLMs are well-aligned with CPU-centric systems, the predominant architecture in such environments, ensuring that \text{BitNet b1.58} can be run efficiently to further boost these devices’ functionalities.

\textbf{New Hardware for 1-bit LLMs}

Emerging efforts, such as Groq\footnote{\url{https://groq.com/}}, highlight the growing success and potential of developing specialized hardware architectures, like LPUs, tailored for LLM workloads. Building on this momentum, we advocate for the development and system-level design of new hardware optimized explicitly for the requirements of 1-bit LLMs, leveraging the computational advances introduced in BitNet~\cite{bitnet}.